\section{Related Work}
\label{sec:related}

\paragraph{Pruning after training.}
Removing weights by loss curvature \cite{lecun1989optimal,hassibi1992second} or by magnitude \cite{han2015learning} showed early on that trained networks are highly redundant. Later work refined the saliency estimate with first-order Taylor expansions \cite{molchanov2017pruning,molchanov2019importance}, empirical Fisher information \cite{theis2018faster} and layer-wise second-order reconstruction \cite{frantar2022obc}. These criteria estimate the increase in loss caused by removing a weight, which is the quantity pruning ultimately needs to control. Large-scale comparisons \cite{gale2019state,blalock2020state,hoefler2021sparsity} found that careful magnitude pruning remains a strong baseline and that many reported gains disappear under matched training budgets, a finding that motivates the controls in this paper.

\paragraph{Pruning at initialisation and gradual pruning.}
SNIP \cite{lee2018snip}, GraSP \cite{wang2020picking} and SynFlow \cite{tanaka2020synflow} fix a mask before training from connection sensitivity or gradient flow; random masks with appropriate layer-wise densities are competitive with them \cite{liu2022random}. Gradual magnitude pruning (GMP) \cite{zhu2017prune} increases sparsity along a cubic schedule during training. Movement pruning \cite{sanh2020movement} keeps weights that move away from zero during training, and dynamic pruning with feedback \cite{lin2020dpf} lets pruned weights recover. The lottery-ticket line of work \cite{frankle2018the,renda2020comparing} studies which sparse sub-networks train well from (near) initialisation.

\paragraph{Dynamic sparse training.}
DST keeps the parameter count fixed throughout training and periodically rewires the network. SET \cite{mocanu2018scalable} drops the smallest-magnitude weights and regrows at random, DeepR \cite{bellec2017deep} rewires stochastically, and DSR \cite{mostafa2019parameter} reallocates parameters across layers. SNFS \cite{dettmers2019sparse} regrows by momentum. RigL \cite{evci2020rigging} regrows by instantaneous gradient magnitude, uses an Erd\H{o}s--R\'enyi-kernel (ERK) layer allocation and decays the update fraction with a cosine schedule; it is the reference DST method and our closest baseline. Subsequent work studied why DST works, attributing it to parameter exploration over time \cite{liu2021itop} and to improved gradient flow \cite{evci2022gradient}. Other work combined DST with gradual pruning \cite{liu2021granet}, kept backward passes sparse as well \cite{jayakumar2020topkast}, targeted memory-limited edge training \cite{yuan2021mest} or learned layer-wise thresholds \cite{kusupati2020soft}. Nearly all of these methods drop by magnitude and change the \emph{regrowth} rule or the schedule. DEGST changes the \emph{drop} rule and keeps RigL's gradient regrowth, so our ablation can attribute any difference to the drop criterion alone.

\paragraph{Structured pruning and activation-based channel selection.}
Structured methods remove whole filters or channels so that the result is a smaller dense network. They score channels by filter norms \cite{li2016pruning}, BatchNorm scaling factors trained under an $\ell_1$ penalty (Network Slimming) \cite{liu2017learning}, distance to the geometric median of a layer's filters (FPGM) \cite{he2019filter}, reconstruction error of the next layer \cite{luo2017thinet,he2017channel}, importance propagated from the output \cite{yu2018nisp}, or learned channel gates (AutoPruner) \cite{luo2020autopruner}. Soft filter pruning \cite{he2018soft} lets pruned filters recover during training, and group-sparsity regularisation \cite{wen2016learning} learns structure directly. Several methods score channels with \emph{activation} statistics: the average percentage of zeros \cite{hu2016network}, the entropy of globally pooled channel responses \cite{luo2017entropy}, the rank of feature maps \cite{lin2020hrank} and the independence between channels \cite{sui2021chip}. DEGST-S uses a per-channel histogram entropy in the same spirit, but inside a prune-and-regrow loop rather than a one-shot selection. In residual and depthwise architectures, channels are coupled across layers; we use DepGraph \cite{fang2023depgraph} to find the coupled groups and to remove them physically.

\paragraph{Information-theoretic views of representations.}
The information bottleneck view of deep learning \cite{tishby2015deep,shwartz2017opening} motivates entropy as a descriptor of how much a representation varies with the input. A caveat applies to all entropy-based criteria, including ours: a channel's marginal entropy measures the diversity of its values, not the sensitivity of the loss to removing the weights that read it. Low entropy suggests redundancy but does not establish it, which is why we compare entropy-based scores directly with loss-sensitivity scores (Taylor and Fisher) in the same loop.

\paragraph{Hardware efficiency of sparsity.}
Irregular sparsity is difficult to accelerate on GPUs. Speed-ups need specialised sparse kernels at high sparsity \cite{gale2020sparse,elsen2020fast} or constrained patterns such as 2:4 semi-structured sparsity supported by recent tensor cores \cite{mishra2021accelerating,zhou2021nm}. Structured pruning reduces latency on standard dense kernels. We therefore report non-zero FLOPs, executed FLOPs and measured latency separately.
